\begin{table}[h!]
  \centering
  \label{tab:types:integer_ranges}
  %% Columns sized to their contents rather than to fixed fractions of the text
  %% width. As p{} columns, the Type cells were too narrow for their two chips
  %% and the Signed Range cells too narrow for their formula, so both wrapped --
  %% which put each type one line below the row it labels, and broke $2^{N-1}-1$
  %% across two lines after the minus sign.
  \begin{tabular}{l|l|c|c}
    \toprule[0.6pt]
    \textbf{Type} & \textbf{Bits} & \textbf{Signed Range} & \textbf{Unsigned Range} \\
    \toprule[0.6pt]
    \token{i8} / \token{u8}      & 8    & $-2^7 \ldots 2^7-1$        & $0 \ldots 2^8-1$ \\
    \token{i16} / \token{u16}    & 16   & $-2^{15} \ldots 2^{15}-1$  & $0 \ldots 2^{16}-1$ \\
    \token{i32} / \token{u32}    & 32   & $-2^{31} \ldots 2^{31}-1$  & $0 \ldots 2^{32}-1$ \\
    \token{i64} / \token{u64}    & 64   & $-2^{63} \ldots 2^{63}-1$  & $0 \ldots 2^{64}-1$ \\
    \token{isize} / \token{usize} & system-dependent (N) & $-2^{N-1} \ldots 2^{N-1}-1$ & $0 \ldots 2^N-1$ \\
    \midrule[0.2pt]
  \end{tabular}
  \caption{Integer Types and Ranges in Ymir (expressed as powers of 2)}
\end{table}
